%%% Load options early.
\PassOptionsToPackage{svgnames}{xcolor}
\PassOptionsToPackage{usenames,dvipsnames}{color}
%%% Maybe it makes sense to use memoir class instead, but this needs some serious adjustments to the template.
\documentclass[10pt, twoside, openright]{book}

%%%%% Packages required for the template. %%%%%

%%% Swedish G5 thesis format
\usepackage[papersize={169mm,239mm},top=10mm, bottom=10mm, textwidth=125mm, includehead, includefoot, marginparwidth = 2.5cm]{geometry}

%%% Set up font.
\usepackage[T1]{fontenc}
\usepackage{plex-sans}

%%% Generally useful typesetting packages.
\usepackage{mathtools,amssymb,amsfonts,amsthm}
\usepackage{newpxtext}
\usepackage{newpxmath}
\usepackage[sectionbib]{chapterbib}
\usepackage{xspace}
\usepackage{fancyhdr}
\usepackage{pageslts}
\usepackage{moresize}
\usepackage[fit]{truncate}
\usepackage{epigraph}
\setlength{\epigraphwidth}{0.9\textwidth}
\setlength{\epigraphrule}{0pt}
\usepackage{appendix}
\usepackage{afterpage}
\usepackage{microtype}
\usepackage{booktabs}
\usepackage[all]{nowidow}
\usepackage[hyphens]{url}
\usepackage[hidelinks]{hyperref}
\usepackage[textfont ={it},labelfont={bf, sf}]{caption}
\usepackage[svgnames]{xcolor}
\usepackage[capitalise, noabbrev]{cleveref}
\newcommand{\creflastconjunction}{, and\nobreakspace}
\usepackage{tocloft}
\usepackage{sectsty}
\usepackage{float}


%%% todo notes setup
\usepackage{todonotes}
%%% To disable todo notes, replace the line \usepackage{todonotes} above by the line below
% \usepackage[disable]{todonotes}


%%% You can include your own packages down here.


\let\oldsection\section% Store \section
\renewcommand{\section}{% Update \section
  \renewcommand{\theequation}{\arabic{equation}}% Update equation number
  \oldsection}% Regular \section
\let\oldsubsection\subsection% Store \subsection
\renewcommand{\subsection}{% Update \subsection
  \renewcommand{\theequation}{\arabic{equation}}% Update equation number
  \oldsubsection}% Regular \subsection


%%% Thesis macros for template %%%
\counterwithin{chapter}{part}
% The macros.tex file contains settings about the layout of the
% Thesis itself. Do not change anything below this line unless you know what
% you are doing.

% ------------------ Change Chapter title row layout
\makeatletter
\def\thickhrulefill{\leavevmode \leaders \hrule height 1ex \hfill \kern \z@}
\def\@makechapterhead#1{%
  \markmargin%
  \vspace*{10\p@}%
  {\parindent \z@ \raggedleft \reset@font
            \sffamily \bfseries %\scshape %\@chapapp{} \thechapter%
        \par\nobreak%
        \interlinepenalty\@M%
    \Huge #1\par\nobreak
    \hrulefill
    \par\nobreak
    \vskip 20\p@
  }}
\def\@makeschapterhead#1{%
  \markmarginKappa%
  \vspace*{10\p@}%
  {\parindent \z@ \raggedleft \reset@font
            \sffamily \bfseries %\scshape %\@chapapp{} \thechapter%
        \par\nobreak
        \interlinepenalty\@M
    \Huge  #1\par\nobreak
    \hrulefill
    \par\nobreak
    \vskip 20\p@
  }}

%------------------------ part page layout
%% from book.cls:
\renewcommand\part{%
   \if@openright
      \cleardoublepage
   \else
      \clearpage
   \fi
   \thispagestyle{plain}%
  \if@twocolumn
    \onecolumn
    \@tempswatrue
  \else
    \@tempswafalse
  \fi
  \null\vfil
  \secdef\@part\@spart}

\def\@part[#1]#2{%
    \ifnum \c@secnumdepth >-2\relax
      \refstepcounter{part}%
      \addcontentsline{toc}{part}{#1}%
    \else
      \addcontentsline{toc}{part}{#1}%
    \fi
    \markboth{}{}%
    {\centering
     \interlinepenalty \@M
     \normalfont
     \Huge \reset@font \hrulefill \\ \raggedleft \vspace{5mm} \sffamily%
            \bfseries {#2} \hspace{15mm}\ \\ \hrulefill \par}%
    \@endpart}
\def\@spart#1{%
    {\centering
     \interlinepenalty \@M
     \normalfont
     \Huge \sffamily \bfseries scshape #1\par}%
    \@endpart}
\def\@endpart{\vfil    %% \newpage -- HWL: no pagebreak with part env
              \if@twoside
                \null
                \thispagestyle{empty}%
                %%\newpage
              \fi
              \if@tempswa
                \twocolumn
              \fi}
\makeatother

\numberwithin{equation}{chapter}
\newcounter{rtaskno}
\newcommand{\paper}[1]{\refstepcounter{rtaskno}\label{#1}}
\renewcommand{\thertaskno}{\Roman{rtaskno}}
\crefname{rtaskno}{Paper}{Papers}
\Crefname{rtaskno}{Paper}{Papers}
%-----------------------------------------------



%Make fancy header and no footer
\pagestyle{fancy}
\renewcommand{\chaptermark}[1]{\markboth{#1}{}}
% \renewcommand{\sectionmark}[1]{\markright{\thesection\ #1}}
\fancyhf{}
\fancyhead[LE, RO]{\bfseries\thepage}
\fancyhead[LO]{\itshape \nouppercase{\rightmark}}
\fancyhead[RE]{\truncate{.95\headwidth}{\itshape \nouppercase{\leftmark}}}
\renewcommand{\headrulewidth}{0.5pt}
\renewcommand{\footrulewidth}{0pt}%
%\setlength{\headheight}{75pt}
\fancypagestyle{plain}{
  \fancyfoot[LE, RO]{\chapterFoot}
  \fancyhead{}
  \renewcommand{\headrulewidth}{0pt}
}

\newcommand{\chapterFoot}{}
\newcommand{\markmargin}{}
\newcommand{\markmarginKappa}{}
\newlength{\spiff}

% -------------------------> Applying the style to the thesis. See macros.tex for more details
\allsectionsfont{\sffamily} % -----------> All sections in sans serif



\newcommand{\Chapter}[1]{Chapter~\ref{#1}}


\newfloat{paperfoot}{b}{paper}
\newcommand{\paperRemark}[1]{
  \begin{paperfoot}%
  \hrulefill \flushleft \footnotesize #1 \end{paperfoot}
}

\renewcommand{\bibname}{References}

\newcommand{\bibpath}{./}

%%% Theorems %%%
\newtheorem{theorem}{Theorem}[chapter]
\renewcommand{\thetheorem}{\arabic{theorem}}
\newtheorem{lemma}[theorem]{Lemma}
\renewcommand{\thelemma}{\arabic{lemma}}
\newtheorem{corollary}[theorem]{Corollary}
\renewcommand{\thecorollary}{\arabic{corollary}}
\newtheorem{proposition}[theorem]{Proposition}
\renewcommand{\theproposition}{\arabic{proposition}}
\newtheorem{observation}[theorem]{Observation}
\renewcommand{\theobservation}{\arabic{observation}}

\theoremstyle{definition}
\newtheorem{example}{Example}[chapter]
\renewcommand{\theexample}{\arabic{example}}
\newtheorem{definition}{Definition}[chapter]
\renewcommand{\thedefinition}{\arabic{definition}}
\newtheorem{remark}{Remark}[chapter]
\renewcommand{\theremark}{\arabic{remark}}

%%% Macros to deal with publisher templates %%%
% Abstracts
\newenvironment{abstract}{\section*{Abstract}}{}

% Appendix
\appendixtitleon
\appendixtitletocon
\renewcommand{\setthesection}{\Alph{section}}
\renewcommand{\setthesubsection}{\Alph{section}.\arabic{subsection}}


%%% Custom macros %%%
%%% TODO: Define macros for including references to your papers inline.
\newcommand{\paperOne}{%
Author 1, Author 2, Author 3, and Author 4.
``Paper with Interesting Title''.
In \emph{Proceedings of the 30th Conference on Very Interesting Topics (CVIT~'24)}, volume 123 of \emph{Series of Very Interesting Conferences (SVIC)}, pages 123\nobreakdash--155, September 2024.
}
\newcommand{\paperTwo}{%
Author 1, Author 2, Author 3, Author 4, and Author 5.
``Another Paper with an Interesting Title''.
In \emph{Journal on Good Thesis Writing}, volume 12345 of \emph{Lecture Notes on Good Thesis}, pages 1\nobreakdash--22. Publisher X, July 2025.
}
\newcommand{\paperNotIncludedOne}{%
Author 1, Author 2, Author 3, Author 4, Author 5, and Author 6.
``Not So Interesting Paper Title''.
In \emph{Proceedings of the 10th International Conference on Some Topics (CST~'23)}, volume 333 of \emph{Series of Very Interesting Conferences (SVIC)}, pages 555\nobreakdash--600, October 2023.
}


%%% You can define more general macros below here.

%%% Set abstract early to be used in the datasheet.
\newcommand{\myAbstract}{
Write you abstract for your thesis here.
}



%%% Set data of the thesis:
%%% TODO: Fill the following command out.
\newcommand{\titlenameshort}{\textcolor{red}{Define titlenameshort}}
\newcommand{\myMainTitle}{\textcolor{red}{Define myMainTitle}}
\newcommand{\mySubtitle}{\textcolor{red}{Define mySubtitle}}
\newcommand{\titlename}{\myMainTitle\xspace\mySubtitle}
\newcommand{\myName}{\textcolor{red}{Define myName}}
\newcommand{\myEmail}{\textcolor{red}{Define myEmail}}
\newcommand{\myWebsite}{\textcolor{red}{Define myWebsite}}
\newcommand{\myYear}{\textcolor{red}{Define myYear}} % year of defence
\newcommand{\myFormDefenceDate}{\textcolor{red}{Define myFormDefenceDate}} % date of defence
\newcommand{\myFormSignDate}{\textcolor{red}{Define myFormSignDate}} % usually the date of printing
\newcommand{\myTitle}{\titlename}
\newcommand{\myFaculty}{\textcolor{red}{Define myFaculty}} % e.g., Faculty of Engineering
\newcommand{\myDepartment}{\textcolor{red}{Define myDepartment}} % e.g., Department of Computer Science
\newcommand{\myAddress}{\textcolor{red}{Define myAddress}} % Adress of university e.g., Box 118\\SE--221 00 LUND\\Sweden
\newcommand{\myFormKeywords}{\textcolor{red}{Define myFormKeywords}}
\newcommand{\myISBNpdf}{\textcolor{red}{Define myISBNpdf}} % electronic version ISBN
\newcommand{\myISBNprint}{\textcolor{red}{Define myISBNprint}} % printed version ISBN
\newcommand{\myISSN}{\textcolor{red}{Define myISSN}}
\newcommand{\myDissertationNumber}{\textcolor{red}{Define myDissertationNumber}}
\newcommand{\myLuCsDissNumber}{\textcolor{red}{Define myLuCsDissNumber}}
\newcommand{\myPages}{\lastpageref{LastPages}}
\newcommand{\myFundingOrganization}{\textcolor{red}{Define myFundingOrganization}} % Funding information for your PhD position, e.g., WASP
\newcommand{\myDegree}{\textcolor{red}{Define myDegree}} % The degree that you want, e.g., Doctor of Philosophy
\newcommand{\myAdvisors}{\textcolor{red}{Define myAdvisors}} % List of your advisors
\newcommand{\myOpponent}{\textcolor{red}{Define myOpponent}}
\newcommand{\myDefenceAnnouncement}{
  To be presented, with the permission of the
	\myFaculty\xspace of Lund University, for public criticism in {room??} at the \myDepartment\xspace on {date??}, \myYear\xspace at {time??}.
} % You need to fill in date, time, and room for your defence.

% For faster compilation you can select to only recompile specific chapters.
% \includeonly{kappa/main}


\begin{document}

\frontmatter

%------------------------------------
% title page
%------------------------------------
%%% First title page with all informat
\cleardoublepage
\thispagestyle{empty}
~
\vfill
\begin{center}
	{\huge \myMainTitle}
	\\[2mm]
	{\Large \mySubtitle}

	\vfill
	{by \myName}
	\vfill
	% black and white (default):
	\includegraphics[width=0.25\textwidth]{LundUniversityLogo.pdf}

	\vspace{10mm}
	{Thesis for the degree of \myDegree}\\
	\vspace{2mm}
	{Thesis advisors: \myAdvisors}\\
	{Faculty opponent: \myOpponent}\\
	\vspace{1cm}
	{\footnotesize
		\myDefenceAnnouncement
	}
	\\
\end{center}
\vfill

%%% Datasheet page
\newpage
\thispagestyle{empty}
\newcounter{aformdx}
\newcounter{aformdy}
\newcounter{aformx}
\newcounter{aformy}
\newcounter{aformtx}
\newcounter{aformtdx}
\newcounter{aformty}
\def\aformw{358}
\def\aformwhalf{\aformw / 2}
\def\aformwquart{\aformw / 4}
\def\aformwthreeq{\aformw / 4 * 3}
\def\aformh{517}
\def\aformb{23}
\def\aformlinew{0.4pt}
\def\aformrowh{20}
\def\aformRowh{24}
\newcommand\aformhline
{\put(\value{aformx},\value{aformy}){\line(1,0){\value{aformdx}}}}
\newcommand\aformvline
{\put(\value{aformx},\value{aformy}){\line(0,-1){\value{aformdy}}}}
\newcommand\aformvlinec
{\put(\value{aformx},\value{aformy}){\line(0,-1){16.25}}}
\newcommand\aformboxbase[5]
{
  \setcounter{aformtx}{\value{aformx}}
  \setcounter{aformtdx}{\value{aformdx}}
  \setcounter{aformty}{\value{aformy}}
  \addtocounter{aformtx}{#1}
  \addtocounter{aformtdx}{-#2}
  \addtocounter{aformty}{-#3}
  \put(\value{aformtx},\value{aformty})
  {
    \parbox[t]{\value{aformtdx}pt}
    {{\sf\tiny #4}{\scriptsize #5\par}}
  }
}
\newcommand\aformbox[2]{\aformboxbase{2}{11}{8}{#1}{#2}}
\newcommand\aformBox[2]{\aformboxbase{2}{11}{8}{#1}{#2}}

\begin{center}
\begin{picture}(\aformw,542)(0,0)
\linethickness{\aformlinew}
\put(0,\aformb){\framebox(\aformw,\aformh){}}
\put(-14,200){\rotatebox{90}{\parbox{200pt}
{\sf\tiny
DOKUMENTDATABLAD
enl SIS 61 41 21\\
}}}
\setcounter{aformx}{0}
\setcounter{aformy}{\aformb}
\addtocounter{aformy}{\aformh}
\setcounter{aformdx}{\aformwhalf}
\aformbox{Organization}
{\\
{\bf LUND UNIVERSITY}\vspace{0.1cm}
\\
{\myDepartment\\
\myAddress}
}
\setcounter{aformdy}{\aformrowh}
\addtocounter{aformdy}{\aformrowh}
\addtocounter{aformdy}{\aformrowh}
\addtocounter{aformdy}{\aformRowh}
\setcounter{aformx}{\aformwhalf}
\aformvline
\setcounter{aformx}{0}
\addtocounter{aformy}{-\aformrowh}
\addtocounter{aformy}{-\aformrowh}
\addtocounter{aformy}{-\aformRowh}
\aformbox{Author(s)}{\\\myName}
\setcounter{aformdy}{\aformrowh}
\aformhline
\addtocounter{aformy}{\aformrowh}
\addtocounter{aformy}{\aformrowh}
\addtocounter{aformy}{\aformRowh}
\setcounter{aformx}{\aformwhalf}
\aformbox{Document name}{\\{\bf DOCTORAL DISSERTATION}}
\addtocounter{aformy}{-\aformrowh}
\aformhline
\aformbox{Date of disputation}{\\\myFormDefenceDate}
\setcounter{aformx}{\aformwhalf}
\aformhline
\addtocounter{aformy}{-\aformrowh}
\addtocounter{aformdy}{\aformRowh}
\aformbox{Sponsoring organization}{\ifdef{\myFundingOrganization}{\\\myFundingOrganization}{}}
\aformhline
\addtocounter{aformy}{-\aformRowh}
\addtocounter{aformy}{-\aformrowh}
\setcounter{aformx}{0}
\setcounter{aformdx}{\aformw}
\aformhline
\aformbox{Title and subtitle }{\\\myTitle}
\addtocounter{aformy}{-\aformRowh}
\setcounter{aformx}{0}
\aformhline
\aformBox{Abstract}
{\\
{\parskip0pt\parindent1em
\myAbstract
}
}
\setcounter{aformy}{\aformb}
% Correspond to one minus each in the bottom section
\addtocounter{aformy}{\aformRowh}
\addtocounter{aformy}{\aformRowh}
\addtocounter{aformy}{\aformRowh}
\addtocounter{aformy}{\aformRowh}
\addtocounter{aformy}{\aformrowh}
\addtocounter{aformy}{\aformrowh}
\addtocounter{aformy}{10}
\addtocounter{aformy}{10}
\aformhline
\aformbox{Key words}
{
\\\myFormKeywords
}
\addtocounter{aformy}{-\aformRowh}
\addtocounter{aformy}{-10}
\aformhline
\aformbox{Classification system and/or index terms (if any)}{}
\addtocounter{aformy}{-\aformRowh}
\aformhline
\aformbox{Supplementary bibliographical information}{}
\setcounter{aformx}{\aformwthreeq}
\setcounter{aformdx}{\aformwquart}
\setcounter{aformdy}{\aformRowh}
\aformvline
\aformbox{Language}{\\English}
\addtocounter{aformy}{-\aformRowh}
\setcounter{aformx}{0}
\setcounter{aformdx}{\aformw}
\aformhline
\aformbox{ISSN and key title}{\\\myISSN}
\setcounter{aformx}{\aformwthreeq}
\addtocounter{aformdy}{10}
\aformvline
\addtocounter{aformdy}{-10}
\aformbox{ISBN}{\\\myISBNprint\xspace(print)\\\myISBNpdf\xspace(pdf)}
\addtocounter{aformy}{-\aformRowh}
\addtocounter{aformy}{-10}
\setcounter{aformx}{0}
\aformhline
\aformbox{Recipient's notes}{}
\setcounter{aformx}{\aformwhalf}
\setcounter{aformdy}{\aformrowh}
\addtocounter{aformdy}{\aformrowh}
\aformvline
\aformbox{Number of pages}{\\\myPages}
\setcounter{aformx}{\aformwthreeq}
\aformbox{Price}{}
\setcounter{aformdy}{\aformrowh}
\aformvline
\addtocounter{aformy}{-\aformrowh}
\setcounter{aformx}{\aformwhalf}
\setcounter{aformdx}{\aformwhalf}
\aformhline
\aformbox{Security classification}{}
\setcounter{aformx}{0}
\setcounter{aformdx}{\aformw}
\addtocounter{aformy}{-\aformrowh}
\addtocounter{aformy}{-4}
\aformbox{}
{\scriptsize I, the undersigned, being the copyright owner of the abstract of the
above-mentioned dissertation, hereby grant to all reference sources
the permission to publish and disseminate the abstract of the
above-mentioned dissertation.
}
\addtocounter{aformy}{-\aformRowh}
\addtocounter{aformy}{-20}
\aformbox{
\hspace{0pt}
%\put(44,1){\includegraphics{images/signature}} %Media-Tryck does not want this to be included as a graphic
}{}
\aformbox{\rm
\scriptsize
Signature \underline{\hskip 140pt}\hfill Date \underline{\hskip 80pt}
}{}
\addtocounter{aformy}{1}
\aformbox{}{\hskip 280pt \scriptsize \myFormSignDate}
\end{picture}
\end{center}


%%% Second title page without additional information to have book metadata on the left-hand side.
\cleardoublepage
\thispagestyle{empty}
~
\vfill
\begin{center}
	{\huge \myMainTitle}
	\\[2mm]
	{\Large \mySubtitle}

	\vfill
	{by \myName}
	\vfill
	% black and white (default):
	\includegraphics[width=0.25\textwidth]{LundUniversityLogo.pdf}

	\color{white}{
	\vspace{10mm}
	{\myDegree}\\
	{Thesis advisors: \myAdvisors}\\
	{Faculty opponent: \myOpponent}\\
	\vspace{1cm}
	{\footnotesize
		\myDefenceAnnouncement
	}}
	\\
\end{center}
\vfill

%%% Book metadata page.
\newpage
\begin{flushleft}
	A doctoral thesis at a university in Sweden takes either the form of a single,
	cohesive research study (monograph) or a summary of research papers
	(compilation thesis), which the doctoral student has written alone or together
	with one or several other author(s).

	In the latter case the thesis consists of two parts. An introductory text puts
	the research work into context and summarizes the main points of the papers.
	Then, the research publications themselves are reproduced, together
	with a description of the individual contributions of the authors. The
	research papers may either have been already published or are manuscripts at
	various stages (in press, submitted, or in draft).

	\vspace*{\stretch{1}}

	\ifdef{\myFundingOrganization}{
		\textbf{Funding information:} \myName\xspace was funded by the \myFundingOrganization.
		\linebreak[2]
	}{}

	ISBN \myISBNpdf{} (electronic version)\\
	ISBN \myISBNprint{} (print version)\\
	ISSN \myISSN{}\\
	Doctoral Dissertation \myDissertationNumber, \myYear\\
	LU-CS-DISS: \myYear-\myLuCsDissNumber\linebreak[2]

	Lund University \\
	\myDepartment \\
	\myAddress  \linebreak[2]

	E-mail: \href{mailto:\myEmail}{\texttt{\myEmail}}\\
	WWW: \href{https://\myWebsite}{\texttt{\myWebsite}} \linebreak[2]

	Typeset using \LaTeX \\

	Printed in Sweden by Tryckeriet i E-huset, Lund, \myYear\linebreak[2]

	\textit{\copyright~\myName{}~\myYear{}}
	\vspace{10mm}
\end{flushleft}

%------------------------------------



% -----------------------------------
% ------- GENERAL STATEMENTS --------
% -----------------------------------
\setcounter{secnumdepth}{3}
\setcounter{tocdepth}{3}
\cftsetindents{chapter}{0em}{2.5em}
{
	\hypersetup{linkcolor=black}
	\tableofcontents
	% \listoffigures
	% \listoftables
}

\chapter{Acknowledgements}
\label{chap:acknowledgements}

I would like to take this opportunity to thank all the people that are important to me.

I would also like to thank the funding organization XYZ for funding my research during my PhD.

\medskip

\hfill \emph{Thank you very much!}

\medskip

\hfill \emph{Tack så mycket!}

\chapter{Popular Science Summary}
\label{chap:pop-sci-sum}

This is the popular science summary of the thesis.
\chapter{Contribution Statement}


The following papers are included in this thesis:

\begin{description}
	\item[\cref{paper:one}]
		\paperOne
	\item[\cref{paper:two}]
		\paperTwo
\end{description}

There are additional papers that \myName\xspace contributed to, which are not included in this thesis:

\begin{itemize}
	\item \paperNotIncludedOne
\end{itemize}

For the included papers, the table below indicates the work preformed by \myName\xspace in the corresponding projects. There were also continuous discussions between all collaborators throughout the project and especially in early stages.

\vspace{1em}
%%% Define contribution circles.
\newcommand{\rbhead}[1]{\rotatebox{70}{\emph{#1}}}
\newcommand*{\cfull}{\tikz[baseline=-3pt]{\fill[black] circle(1ex);}}
\newcommand*{\cempt}{\tikz[baseline=-3pt]{\draw circle(1ex);}}
\newcommand*{\cquarter}{\tikz[baseline=-3pt]{\fill[black] (0,0) -- (0:1ex) arc (0:90:1ex) -- cycle;%
		\draw (0,0) -- (0:1ex) arc (0:360:1ex);}}
\newcommand*{\chalf}{\tikz[baseline=-3pt]{\fill[black] (0,0) -- (90:1ex) arc (90:-90:1ex) -- cycle;%
		\draw (0,0) circle (1ex);}}
\newcommand*{\cthreequarter}{\tikz[baseline=-3pt]{\fill[black] (0,0) -- (90:1ex) arc (90:-180:1ex) -- cycle;%
		\draw (0,0) -- (0:1ex) arc (0:360:1ex);}}
\begin{center}
	\begin{tabular}{lcccc}
		\toprule
		\emph{Paper}                           & \emph{Concept} & \emph{Implementation} & \emph{Evaluation} & \emph{Writing}                  \\
		\midrule
		\textbf{\ref{paper:one}}         & \chalf & \cthreequarter    & \cquarter      & \chalf         \\ % First paper
		\textbf{\ref{paper:two}}       & \cthreequarter & \cthreequarter    & \cfull         & \cthreequarter \\ % Second paper
		\bottomrule
	\end{tabular}
\end{center}
\vspace{1em}

\noindent
The dark portion of the circle is intended to illustrate the proportion of work and responsibilities assigned to \myName\xspace, but is more formally defined as:
\begin{itemize}
	\item[\cfull] \myName\xspace led and did almost all the work.
	\item[\cthreequarter] \myName\xspace led and did a majority of the work.
	\item[\chalf] \myName\xspace was a contributor to the work.
	\item[\cquarter] \myName\xspace was a minor contributor to the work.
	\item[\cempt] \myName\xspace did not contribute, except for proofreading.
\end{itemize}

\medskip

\noindent\textbf{Concept}\enskip Coming up with the ideas and working out the details in theory.
\smallskip

\noindent\textbf{Implementation}\enskip Writing the software required for the paper.
\smallskip

\noindent\textbf{Evaluation}\enskip Conducting the experimental evaluation and analysing the data.
\smallskip

\noindent\textbf{Writing}\enskip Drafting and editing the paper.


\medskip

\noindent The following discusses the contributions by each author and especially the contribution by \myName\xspace in more detail.


% Paper 1
\subsection*{\cref{paper:one}}
Author 1 did this. Author 2 did that. \dots

% Paper 2
\subsection*{\cref{paper:two}}
Author 2 and Author 4 did that. Author 5 did that. All authors were involved in something. \dots



\mainmatter
\renewcommand{\thesection}{\arabic{section}}
\renewcommand{\thefigure}{\arabic{figure}}
\renewcommand{\thetable}{\arabic{table}}
\renewcommand{\thechapter}{\Roman{chapter}}
\renewcommand{\theequation}{\arabic{equation}}

% Add right margin marks for chapters.
%-------------------------------------------------------
\newlength{\apa}
\setlength{\spiff}{0cm}
\renewcommand{\markmargin}{%
	\vspace*{\spiff}
	\marginpar{\raggedright \rule{5mm}{0mm}
		\begin{tikzpicture}
			\node[rectangle, fill=black, minimum width=6cm, minimum height=2cm, inner sep=0pt,draw=black](rec){};
			\node[left of = rec,node distance=2.7cm,text width=2cm,text centered, font=\sffamily\bfseries\small,rotate=90]{\color{white}{Paper \thechapter}};
		\end{tikzpicture}\rule{1pt}{\apa}}%
	\vspace*{-\spiff}
	\addtolength{\spiff}{2cm}
}
\renewcommand{\markmarginKappa}{%
	\vspace*{\spiff}
	\marginpar{\raggedright \rule{5mm}{0mm}
		\begin{tikzpicture}
			\node[rectangle, fill=black, minimum width=6cm, minimum height=2cm, inner sep=0pt,draw=black](rec){};
			\node[left of = rec,node distance=2.7cm,text width=2cm,text centered, font=\sffamily\bfseries\small,rotate=90]{\color{white}{Kappa}};
		\end{tikzpicture}\rule{1pt}{\apa}}%
	\vspace*{-\spiff}
	\addtolength{\spiff}{2cm}
}
%-------------------------------------------------------

% --------------------------------
% --------- KAPPA ----------------
% --------------------------------
\chapter*{\texorpdfstring{\titlename\hspace{-0.2cm}}{\titlename}}
\addcontentsline{toc}{part}{\titlenameshort}
\markboth{\titlenameshort}{}

% Introduction
\section{Introduction}
\label{sec:kappa:introduction}

This is the introduction of the kappa. This work is important~\cite{Alon00Algebraic,Alon99Combinatorial,AR23MCSP}. This problem has been studied in the literature~\cite{AN05Chromatic,AR22PerfectMatchingJournal,BBKMW21SpectralPlanting}.

The thesis is structured as follows. First, the so-called kappa gives a comprehensive summary over the thesis. After that the papers included in this thesis are given. The kappa starts with some preliminaries in \cref{sec:kappa:preliminaries}. \dots

% Preliminaries/Background
\section{Background}
\label{sec:kappa:preliminaries}

Here you can write about background and preliminaries for your thesis.

% Related work
\section{Related Work}
\label{sec:kappa:related}

Some work related to the thesis is mentioned here.

% Contribution
\section{Main Results of Included Papers}
\label{sec:contributions}

In this section an overview of the papers included in this thesis is given to highlight the contribution of each paper included in this thesis to the field.
While \cref{paper:one} focuses on foo, \cref{paper:two} focuses on bar.

\subsection{Summary of \texorpdfstring{\cref{paper:one}}{Paper I}}
\epigraph{\paperOne}{}

\noindent
Summarize the contributions of your first paper here.


\subsection{Summary of \texorpdfstring{\cref{paper:two}}{Paper II}}
\epigraph{\paperTwo}{}

\noindent
Summarize the contributions of your second paper here.


\subsection{Further Contributions}
Maybe you have some further contributions that you want to include here.

% Future work and conclusion
\section{Conclusions and Future Work}%
\label{sec:kappa:conclusions}

The conclusions are that.

We leave the following things as interesting future work.

\bibliographystyle{alpha}
\bibliography{refLocal}




\renewcommand{\thechapter}{\Roman{chapter}}
\renewcommand{\chaptername}{Paper}
\setcounter{rtaskno}{0}

% --------------------------------
% --------- PAPERS ---------------
% --------------------------------

\part{Included Papers}
%%% TODO: Include the main files of your included papers here.

%%% TODO: Set the title for the paper.
\chapter[Paper with Interesting Title]{Paper with Interesting Title}

%%% TODO: Add reference for the paper.
\paperRemark{\paperOne}
%%% TODO: Add a label to the paper.
\paper{paper:one}

%%% TODO: Set up paths for included graphics and bib files.
\makeatletter
  \def\input@path{{papers/one/}}
  \graphicspath{{papers/one/}}
\makeatother
\renewcommand{\bibpath}[1]{papers/one/#1}

%%% Add paper specific macros.
\input{macros-local}

%%% Input abstract of paper.
\input{abstract}

%%% TODO: Files of the paper.
\input{introduction}
\input{preliminaries}
\input{contributions}
\input{conclusion}

\input{acknowledgements}

%%% TODO: Add bib for paper.
\bibliographystyle{alpha}
\bibliography{\bibpath{refLocal}}


%%% TODO: Set the title for the paper.
\chapter[Another Paper with an Interesting Title]{Another Paper with an Interesting Title}

%%% TODO: Add reference for the paper.
\paperRemark{\paperTwo}
%%% TODO: Add a label to the paper.
\paper{paper:two}

%%% TODO: Set up paths for included graphics and bib files.
\makeatletter
  \def\input@path{{papers/two/}}
  \graphicspath{{papers/two/}}
\makeatother
\renewcommand{\bibpath}[1]{papers/two/#1}

%%% Add paper specific macros.
\input{notationLocal.tex}

%%% Input abstract of paper.
\input{abstract.tex}

%%% TODO: Files of the paper.
\input{introduction.tex}
\input{preliminaries.tex}
\input{background.tex}
\input{concludingremarks.tex}

%%% You can also add appendices for the paper.
\begin{subappendices}
  \input{appendix.tex}
\end{subappendices}


%%% TODO: Add bib for paper.
\bibliographystyle{alpha}
\bibliography{\bibpath{refLocal}}


\end{document}
